% Generated draft: requires mathematical and source review.
\documentclass[11pt]{article}
\usepackage[margin=25mm]{geometry}
\usepackage{fontspec}
\setmainfont{Latin Modern Roman}
\usepackage{amsmath,amssymb,mathtools,unicode-math}
\setmathfont{Latin Modern Math}
\usepackage{enumitem,needspace,xurl,hyperref}
\setlength{\parindent}{0pt}
\setlength{\parskip}{5pt}
\setlength{\emergencystretch}{3em}
\newcommand{\N}{\mathbb{N}}
\newcommand{\Z}{\mathbb{Z}}
\newcommand{\Q}{\mathbb{Q}}
\newcommand{\R}{\mathbb{R}}
\newcommand{\C}{\mathbb{C}}
\newcommand{\F}{\mathbb{F}}
\newcommand{\A}{\mathbb{A}}
\newcommand{\PP}{\mathbb{P}}
\newcommand{\E}{\mathbb{E}}
\newcommand{\eps}{\varepsilon}
\newcommand{\dd}{\,\mathrm{d}}
\begin{document}
\section*{UnsolvedMath ID: 700028}
\textbf{Source problem code:} catalogue entry
\subsection{Definitions and mathematical statement}
For a set \(X\), let \(F(X)\) be the free group on \(X\). Define its derived series by
\[
F(X)^{(0)}=F(X),\qquad F(X)^{(i+1)}=[F(X)^{(i)},F(X)^{(i)}],
\]
where \([H,H]\) is the subgroup generated by the commutators \(a^{-1}b^{-1}ab\) with \(a,b\in H\). For a positive integer \(d\), the standard free solvable group on \(X\) of derived length at most \(d\) is
\[
S(X,d)=F(X)/F(X)^{(d)}.
\]
Its automorphism group \(\operatorname{Aut}(S(X,d))\) consists of all bijective group homomorphisms from \(S(X,d)\) to itself, with composition as the group operation.

\textbf{Problem.} Describe the automorphism group of a free solvable group; in the standard notation above, describe
\[
\operatorname{Aut}(S(X,d)).
\]
The source specifies neither the rank \(|X|\) nor the derived-length parameter \(d\), and gives no restriction to finite rank or nontrivial boundary cases. The notation above makes the conventional family explicit, but the intended parameter range remains unspecified. The source also does not prescribe what constitutes a satisfactory description of the automorphism group.
\subsection{Short English statement}
Describe the group of all automorphisms of a free solvable group. Such a group is conventionally obtained from a free group by factoring out a term of its derived series. The intended rank, derived-length range, and form of the requested description are unspecified.
\subsection{Sources}
\begin{itemize}
\item Local supplied UnsolvedMath registry record: permanent ID \(700028\), source problem code \texttt{catalogue entry}, title \emph{Kourovka Notebook Problem 1.33}; upstream revision \texttt{d371ee917ac9647fbb06cd6fc529deec6f3dda82}.
\item Supplied bibliographic source: \emph{The Kourovka Notebook (Unsolved Problems in Group Theory)}, original problem 1.33, first issue (1965), as represented in the twenty-first issue edition (2026), arXiv:1401.0300v48. The supplied collection description names E. I. Khukhro and V. D. Mazurov as editors.
\item Attribution: the record names M. I. Kargapolov as proposer and gives 1965 as the proposal year; the problem statement itself begins with the attribution ``(A. I. Mal'cev).'' Both attributions are retained without resolving their respective roles.
\item The registry labels the problem \texttt{open}. This is supplied source metadata, not an independently verified assessment of its mathematical or literature status.
\end{itemize}
\end{document}
